\section{Conclusions}
In this paper, we have measured an escape velocity profile of the Milky Way from $4-11\,\rm{kpc}$ using a sample of high-speed stars from \Gaia~DR3 with 6D kinematics and strict quality cuts. Following \cite{Necib2022methods}, we considered single and multi-component power laws to model the tail of the stellar speed distribution and extract escape velocities. Similar to \cite{Necib2022}, we found that a single power law fit does not provide a good fit to the tail and that it often significantly overestimates the escape velocity compared to the fastest star in the data. The multi-component power law model is motivated by the presence of kinematic substructure, and provides a much better fit. However, with the high quality data of \Gaia~DR3 and a large range of radii, we found even the two-component power law tends to systematically predict higher $\vesc$. This is particularly noticeable at positions beyond the solar radius, where the power law models give a rising escape velocity profile. Furthermore the $\vesc$ result is highly correlated with the slope of the distribution at low speeds, which is often an unconverged parameter in the fits.

Motivated by the above issues with the power law models, we introduced an empirical model which we call the ``stretched exponential power law" (SEPL). This model was motivated by the same features that multi-component power laws aimed to capture, such as a different behavior near $\vesc$ and near $\vmin$. However, in this case, there is more freedom for a steep rise in the distribution near $\vmin$, while the profile near $\vesc$ remains that of a power law. The escape velocity profile obtained using this new model is less susceptible to the issue of overestimating $\vesc$, and gives a falling profile past the solar position, consistent with expectation. This model is empirically motivated, and future work may improve this modeling as data quality and statistics improve or by studying its application to simulations.
%, to quantify the relationship between the measured escape velocity and the true underlying $\vesc$.   

Using the resulting escape velocity profile, we constrained the DM halo parameters of the Milky Way, in particular its virial mass, by assuming a model of the baryonic content consistent with the literature. By combining constraints with a complementary circular velocity measurement, we find a total virial mass for the Milky Way  of $M_{200\rm{c}} = 0.64^{+0.15}_{-0.14}\times 10^{12} \, \rm{M}_\odot$ using the SEPL model for speed distributions. We obtain consistent but less constrained estimates using the 2PL model. This Milky Way mass is found to be consistent with recent measurements in the literature using power law models, albeit on the lighter end, following the recent trend with such analyses. This trend may result from the improvement in data as well as more modeling techniques that account for substructure in the speed distribution. It has also been demonstrated in \cite{Kravtsov:2024} that a massive satellite galaxy such as the Large Magellanic Cloud can lead to the development of a high-velocity tail in the host, which would bias mass inferences made using stellar speeds. As a result, future modelling may require more sophisticated Milky Way system mass models, or may utilize comparisons to simulated Milky Way analogs including massive companions \citep{Buch:2024}. 

\begin{figure}[t]
    \includegraphics[width=\linewidth]{media/MW_mass_comparison.pdf}
    \caption{Milky Way total mass estimates obtained via escape velocity modeling along with  circular velocity constraints. We compare results from the literature with those of this work. All error bars correspond to $68\%$ confidence intervals, and arrows imply corrections to the inferred escape velocities motivated by simulations.}
    \label{fig:mw_mass_comp}
\end{figure}

Apart from the escape velocity, the mass of the Milky Way has been inferred by numerous methods, each resulting in estimates distributed roughly between $0.5-2\times 10^{12}\,\rm{M}_\odot$ over the past decade \citep{wang:2020}. Masses inferred via escape velocities have constituted some of the larger estimates in this population in the past; however, as mentioned above, recent estimates trend below $10^{12}\,\rm{M}_\odot$. Interestingly, some measurements of masses inferred via other probes such as the rotation curve of the inner Milky Way \citep{Ou:2023, Francesco:2023, 2023arXiv230900048J, 2023ApJ...946...73Z} and comparing the Milky Way satellite population to simulations \citep{barber:2014, cautun:2014, patel:2018} have also trended downward in the past decade. This sits in contrast to some recent Galaxy mass estimates obtained using tracer stars \citep{Shen:2022} and satellite dwarfs \citep{fritz:2020, Slizewski:2022}, which suggest an intermediate-mass Milky Way ($\sim 1.4\times 10^{12}\,\rm{M_\odot}$). Other probes such as phase space distribution modeling \citep{li:2020, callingham:2019} and modeling the orbital dynamics of globular clusters \citep{sohn:2018,watkins:2019,Jianling:2022} exhibit no clear trend. See \cite{wang:2020} and \cite{Sawala:2023} for reviews of the Milky Way mass obtained by various methods. These results together paint a picture of a Milky Way whose mass is close to $10^{12}\,M_\odot$, although with some disagreement to within a factor $\lesssim 2$. As observations continue to rapidly improve, a unified understanding of these different facets of the Milky Way's dynamics will be essential to ultimately pinning down the characteristics and history of our Galaxy.
